\SourceLeaf{295}%
und vom Mittel der beiden Azimute ausgeht. Schreiben wir
\[
\tfrac{1}{2}(S+S')=B,\quad\tfrac{1}{2}(T+T')=A,\quad S-S'=b,\quad T-T'=a
\]
so erhalten wir zunächst die Formeln
\[
\begin{aligned}
\sin\tfrac{1}{2}r\sin A&=\sin\tfrac{1}{2}\lambda\cos B\\
\sin\tfrac{1}{2}r\cos A&=\cos\tfrac{1}{2}\lambda\sin\tfrac{1}{2}b\\
\cos\tfrac{1}{2}r\sin\tfrac{1}{2}a&=\sin\tfrac{1}{2}\lambda\sin B\\
\cos\tfrac{1}{2}r\cos\tfrac{1}{2}a&=\cos\tfrac{1}{2}\lambda\cos\tfrac{1}{2}b
\end{aligned}
\]
Betrachtet man also \(A,B,r\) als gegeben, so bestimmt man \(a\) und \(\lambda\) aus den Formeln
\[
\begin{aligned}
\sin A\operatorname{tang}B\operatorname{tang}\tfrac{1}{2}r&=\sin\tfrac{1}{2}a\\
\frac{\sin A\sin\tfrac{1}{2}r}{\cos B}&=\sin\tfrac{1}{2}\lambda
\end{aligned}
\]
und anschließend \(b\) aus
\[
\frac{\cos A\operatorname{tang}\tfrac{1}{2}r}{\cos\tfrac{1}{2}a}=\operatorname{tang}\tfrac{1}{2}b
\]
oder aus
\[
\frac{\cos A\sin\tfrac{1}{2}r}{\cos\tfrac{1}{2}\lambda}=\sin\tfrac{1}{2}b
\]
Da \(r,a,\lambda,b\) klein sind, wird man jedoch lieber die folgenden Formeln verwenden. Sie sind viel bequemer und bis zur fünften Ordnung ausschließlich genau:
\[
\begin{aligned}
a^0&=r\sin A\operatorname{tang}B\\
\lambda^0&=\frac{r\sin A}{\cos B}\\
b^0&=r\cos A\\
\log a&=\log a^0+\mu rr+\tfrac{1}{2}\mu a^0a^0\\
\log\lambda&=\log\lambda^0-\tfrac{1}{2}\mu rr+\tfrac{1}{2}\mu\lambda^0\lambda^0\\
\log b&=\log b^0+\tfrac{1}{2}\mu a^0a^0+\mu\lambda^0\lambda^0
\end{aligned}
\]
Wie man sieht, ist die dritte Korrektur gleich der Summe aus der ersten und dem Doppelten der zweiten.

Für unsere Aufgabe liefern diese Formeln zwar keine direkte Lösung, doch kann man sie zur Kontrolle oder als knappe, übersichtliche Zusammenfassung der direkten Lösung verwenden. Wer aber etwas Übung im numerischen Rechnen hat, kann sie auch leicht für eine indirekte Lösung nutzen und wird dieser wegen ihrer Bequemlichkeit und Genauigkeit vor allen anderen Lösungen den Vorzug geben, besonders wenn für andere Zwecke schon eine grobe Näherung vorliegt.
\par\medskip\centerline{\rule{.25\textwidth}{.3pt}}
